\chapter*{Scheda del documento}
\addcontentsline{toc}{chapter}{Scheda del documento}
\markboth{Scheda del documento}{}

\begin{table}[h]
\centering
\begin{tabularx}{\linewidth}{@{}lL@{}}
\toprule
\textbf{Progetto}      & \progetto{} -- \sottotitolo \\
\textbf{Documento}     & \titolodoc \\
\textbf{Versione}      & \versionedoc \\
\textbf{Stato}         & \statodoc \\
\textbf{Insegnamento}  & \insegnamento{}, A.A.\ \annoaccademico \\
\textbf{Autori}        & \autoreuno, \autoredue \\
\textbf{Repository}    & \url{https://github.com/terrecruis/kintsugi} \\
\bottomrule
\end{tabularx}
\end{table}

\section*{Storico delle revisioni}

\begin{table}[h]
\centering
\begin{tabularx}{\linewidth}{@{}llL@{}}
\toprule
\textbf{Versione} & \textbf{Data} & \textbf{Modifiche} \\
\midrule
0.1 & 22/10/2026 & Prima bozza: obiettivi, architettura V1--V3, ambienti, modello dei guasti, definizioni di misura. \\
0.2 & 22/10/2026 & Riorganizzazione in forma di documentazione: frontespizio, scheda, acronimi, capitoli, decisioni in formato ADR, appendice di pianificazione. \\
0.3 & 22/10/2026 & Aggiunte V4 (database replicato con failover) e V5 (autoscaling sulla coda) e il guasto F7 (blocco). Revisione critica del piano: metodo di iniezione, timestamp, sonde, deadline, SLO e budget di errore, cluster a tre nodi, minacce alla validità, tabella dei parametri. \\
0.4 & 22/10/2026 & V4 e V5 accettate (estensione lasciata alla scelta del gruppo); cluster kind creato e versione di Kubernetes fissata. \\
0.5 & 22/10/2026 & Documento semplificato: linguaggio meno tecnico, capitoli più brevi, decisioni riassunte. \\
0.6 & 23/10/2026 & Decisioni D0001--D0004: sincrono o con coda, configurazione della coda, riprova dei worker, ordini doppi. \\
0.7 & 23/10/2026 & Limite del control plane unico; sonde standard di F7 e tempi previsti; V3 con 1 e con 10 worker come riferimenti per V5, costo in worker-secondi; esperimenti da 225 a 245. D0013, D0016 e D0017 accettate (computer di misura, iniezione dei guasti, SLO). \\
0.8 & 27/10/2026 & Revisione formale. Reintegrate le definizioni di misura (disponibilità di accettazione e di completamento, ripristino, categorie degli ordini), la struttura del run, i carichi W1--W3 e la statistica. Limiti classificati per tipo di minaccia. Nuovi diagrammi: metodo, ambiente locale e AWS, anatomia di un fermo, linea del tempo di un run, albero delle categorie, matrice degli esperimenti, Gantt e ore per fase. Corretti D0005 (orologio comune), punti unici di guasto di V4s e percorso di Toxiproxy. \\
\bottomrule
\end{tabularx}
\end{table}

\section*{Convenzioni}

\begin{itemize}
  \item Le parti \todo{in rosso} sono ancora da scrivere: il documento cresce insieme al progetto, milestone dopo milestone.
  \item Le decisioni architetturali sono identificate da \textbf{D}\textit{nnnn} (capitolo~\ref{cap:decisioni}); le versioni da \textbf{V1}--\textbf{V5}; le domande di ricerca da \textbf{RQ1}--\textbf{RQ6}; i guasti da \textbf{F1}--\textbf{F7}; i carichi da \textbf{W1}--\textbf{W3}; le milestone da \textbf{M1}--\textbf{M7}.
  \item I termini tecnici sono spiegati alla prima occorrenza e raccolti in \emph{Acronimi e termini}; i nomi dei componenti (API, worker, broker) restano in inglese, come nella documentazione degli strumenti.
  \item Nei diagrammi il colore arancione indica sempre un elemento guastato o un istante misurato.
  \item Il diario delle operazioni, richiesto come consegna separata, è nel repository in \texttt{report/diario/}.
\end{itemize}
